\documentclass[../../../../main.tex]{subfiles}
\begin{document}

\begin{flushright}
\footnotesize\textit{【自動文字起こし・要確認】}
\end{flushright}

\noindent
\textbf{[解]} 題意を M.I. で示す. $n=1$ の時
\[
f'(x) = \frac{1 - \log x}{x^2}
\]
だから $a_1 = 1, b_1 = -1$ とすれば良い. $n=k$ での成立を仮定する.
\begin{align*}
f^{(k+1)}(x) &= \frac{b_k \cdot x^k - (k+1) \cdot x^k (a_k + b_k \log x)}{(x^{k+1})^2} \\
&= \frac{(b_k - (k+1) a_k) + (-(k+1) b_k) \log x}{x^{k+2}}
\end{align*}
だから
\[
\begin{cases}
a_{k+1} = -(k+1) a_k + b_k \\
b_{k+1} = -(k+1) b_k
\end{cases}
\]
とすれば $n=k+1$ でも成立.

\bigskip

\noindent
以上より示された. 同漸化式は
\[
\begin{cases}
a_{n+1} = -(n+1) a_n + b_n &, a_1 = 1 \\
b_{n+1} = -(n+1) b_n &, b_1 = -1
\end{cases}
\]

\bigskip

\noindent
(2) (1)から, $b_n = (-1)^n n!$ であるから,
\[
a_{n+1} = -(n+1) a_n + (-1)^n \cdot n! \quad \cdots \textcircled{1}
\]
$c_n = \frac{a_n}{(-1)^n n!}$ とおいて ($c_1 = -1$), \textcircled{1}の両辺を $(-1)^{n+1} (n+1)!$ で割ると
\[
c_{n+1} = c_n + \frac{1}{-(n+1)}
\]
だから, 階差数列より $n \ge 2$ の時
\begin{align*}
c_n &= c_1 + \sum_{k=2}^n \frac{1}{-k} \\
&= -h_n
\end{align*}
$n=1$ でもこれは成立するから
\[
a_n = (-1)^{n+1} \cdot n! \cdot h_n
\]
である.

\newpage

\end{document}
